\chapter{متتاليات الدوال ومتسلسلاتها}\label{ch:b2:funcseq}

حين تتقارب دوال إلى دالة، أي الخواص ينجو من
الانتقال إلى النهاية؟ \omterm{def:b2:funcseq:def}{التقارب النقطي} لا يحفظ شيئًا
تقريبًا؛ أما التقارب \emph{المنتظم} --- أي التقارب \omterm{def:b2:nvs:norm}{بمعيار} النهاية العليا
--- فيحفظ الاتصال والتكاملات على القطع، ويحفظ المشتقات بلمسة
واحدة. ويبرهن هذا الفصل على مبرهنات النقل الثلاث،
وعلى نسخها للمتسلسلات، ويتوّجها بمبرهنة فايرشتراس
في التقريب، مبرهَنةً \omterm{thm:b2:funcseq:weierstrass}{بكثيرات حدود برنشتاين} الاحتمالية
الجميلة.

\section{التقارب النقطي والمنتظم}

\begin{definition}\label{def:b2:funcseq:def}
لتكن $f_n, f \colon X \to \R$ (أو $\C$، أو فضاءً معياريًا)، ولتكن $X$ مجموعة
كيفية. تتقارب $(f_n)$ إلى $f$
\emph{نقطيًا}\index{تقارب نقطي} إذا كان $f_n(x) \to
f(x)$ من أجل كل $x$؛ و\emph{بانتظام}\index{تقارب منتظم}
إذا كان
\[
\norm{f_n - f}_\infty = \sup_{x \in X}\, \abs{f_n(x) - f(x)}
\xrightarrow[n \to \infty]{} 0 .
\]
والمنتظم يستلزم النقطي؛ وعلى $C(\intcc{a}{b})$، يكون التقارب المنتظم
بالضبط هو التقارب في \omterm{def:b2:nvs:banach}{فضاء باناخ}
$\bigl(C(\intcc{a}{b}), \norm\cdot_\infty\bigr)$ المكوَّن من
\cref{ch:b2:nvs}.
\end{definition}

\begin{example}\label{ex:b2:funcseq:xn}
على $\intcc{0}{1}$، تتقارب $f_n(x) = x^n$ \omterm{def:b2:funcseq:def}{نقطيًا} إلى النهاية
\emph{غير المتصلة} $f = \mathbf{1}_{\{1\}}$؛ والتقارب
ليس منتظمًا: $\norm{f_n - f}_\infty \geq f_n\bigl(1 -
\tfrac1n\bigr) = (1 - \tfrac1n)^n \to \eu^{-1} \neq 0$. وعلى
$\intcc{0}{a}$ مع $a < 1$ \emph{يكون} منتظمًا ($\sup = a^n \to
0$): فالانتظام خاصية لمجال التعريف بقدر ما هو خاصية
للمتتالية.
\end{example}

\begin{omfigure}
\begin{tikzpicture}
\begin{axis}[omaxis, width=8.2cm, height=5.6cm,
    xmin=0, xmax=1.08, ymin=0, ymax=1.15,
    xtick={0,1}, ytick={1}]
  \addplot[omDef!40, thick, domain=0:1, samples=80] {x^2};
  \addplot[omDef!60, thick, domain=0:1, samples=80] {x^5};
  \addplot[omDef, thick, domain=0:1, samples=120] {x^15};
  \addplot[omThm, very thick, domain=0:1, samples=160] {x^40};
  \node[font=\small, omDef] at (axis cs:0.62,0.5) {$x^2$};
  \node[font=\small, omDef] at (axis cs:0.83,0.48) {$x^5$};
  \node[font=\small, omThm] at (axis cs:0.987,0.5) {$x^{40}$};
\end{axis}
\end{tikzpicture}

{\small المتتالية $x^n$ على $\intcc{0}{1}$: تتهاوى المنحنيات نحو
$0$ لكن عليها جميعًا أن تتسلق إلى $1$ عند $x = 1$ --- فمسافة النهاية العليا إلى
النهاية النقطية غير المتصلة لا تنكمش تحت ثابت أبدًا.}
\end{omfigure}

\begin{example}[نهايتان ترفضان التبادل]\label{ex:b2:funcseq:doublelimit}
الفصل كله عن تبادل النهايات، وهذا أصغر
فشل ممكن. لتكن $a_{n,m} = \dfrac{n}{n+m}$ من أجل
$n, m \geq 1$. عندئذٍ
\[
\lim_{m\to\infty}\Bigl(\lim_{n\to\infty}a_{n,m}\Bigr)
= \lim_{m\to\infty} 1 = 1,
\qquad
\lim_{n\to\infty}\Bigl(\lim_{m\to\infty}a_{n,m}\Bigr)
= \lim_{n\to\infty} 0 = 0 :
\]
فالنهايتان المتكررتان موجودتان وهما مختلفتان. وكل مبرهنة نقل
في هذا الفصل رخصةٌ لتبادل نهايتين ---
$\lim_n$ مع $\lim_{x\to a}$ (الاتصال)، ومع $\int$
(المكاملة)، ومع $\frac{\dd}{\dd x}$ (الاشتقاق) ---
\omterm{def:b2:funcseq:def}{والتقارب المنتظم} هو بالضبط الرسم الذي يجعل التبادل
مشروعًا. والفكرة الختامية: كلما بدّل ``برهان''
عمليتَي نهاية صامتًا، فهذا الجدول ذو السطرين هو
المثال المضاد الذي يُرفع في وجهه؛ والنتوءات المنزلقة في
\cref{exo:b2:funcseq:2} هي الظاهرة نفسها بلباس علامة
تكامل.
\end{example}

\section{مبرهنات النقل الثلاث}

\begin{theorem}[الاتصال]\label{thm:b2:funcseq:continuity}
إذا كانت كل $f_n$ متصلة عند $a$ وكان $f_n \to f$ بانتظام على
جوار للنقطة $a$، فإن $f$ متصلة عند $a$. والنهاية المنتظمة
لدوال متصلة متصلة.
\end{theorem}

\begin{proof}
حجة $3\varepsilon$ المستعملة بالفعل في
\cref{thm:b2:metric:rncomplete}: نختار $n$ بحيث $\norm{f_n -
f}_\infty \leq \varepsilon$، ثم $\delta$ من اتصال
$f_n$ عند $a$؛ فمن أجل $\abs{x - a} \leq \delta$،
\[
\abs{f(x) - f(a)} \leq \abs{f(x) - f_n(x)} + \abs{f_n(x) - f_n(a)}
+ \abs{f_n(a) - f(a)} \leq 3\varepsilon . \qedhere
\]
\end{proof}

\begin{example}[يفشل الانتظام بالضبط حيث تنكسر النهاية]\label{ex:b2:funcseq:threshold}
على $\intcc{0}{2}$، لتكن $f_n(x) = \dfrac{x^n}{1 + x^n}$. النهاية
النقطية دالة من ثلاث قطع:
\[
f(x) = \begin{cases} 0 & 0 \leq x < 1,\\[2pt]
\tfrac12 & x = 1,\\[2pt]
1 & 1 < x \leq 2, \end{cases}
\]
وهي غير متصلة عند $1$، ومن ثم لا يمكن أن يكون التقارب
منتظمًا على $\intcc{0}{2}$ حسب \cref{thm:b2:funcseq:continuity}. لكنه منتظم على القطع
المغلقة التي تتجنب العتبة: فمن أجل $0 \leq x \leq
a < 1$،
\[
\sup_{\intcc{0}{a}}\abs{f_n - 0}
= \frac{a^n}{1 + a^n} \leq a^n \to 0 ,
\]
ومن أجل $1 < b \leq x \leq 2$،
\[
\sup_{\intcc{b}{2}}\abs{f_n - 1}
= \frac{1}{1 + b^n} \leq b^{-n} \to 0 ,
\]
وقد حُسبت النهايتان العلييان برتابة $u \mapsto
\frac{u}{1+u}$ ورتابة $x \mapsto x^n$. والفكرة الختامية: يتوطّن
فشل الانتظام عند نقطة عدم اتصال
النهاية --- بالهندسة نفسها التي في \cref{ex:b2:funcseq:xn}، وهو
سبب تكرار انضباط ``الانتظام على كل قطعة في الداخل''
طوال الفصل.
\end{example}

\begin{theorem}[المكاملة على قطعة]\label{thm:b2:funcseq:integration}
إذا كان $f_n \to f$ بانتظام على $\intcc{a}{b}$، مع $f_n$ متصلة
على قطع (و$f$ كذلك)، فإن
\[
\int_a^b f_n \longrightarrow \int_a^b f .
\]
\end{theorem}

\begin{proof}
تعطي الخطية ومتراجحة المثلث للتكاملات أن
\[
\Bigl|\int_a^b f_n - \int_a^b f\Bigr|
= \Bigl|\int_a^b (f_n - f)\Bigr|
\leq \int_a^b\abs{f_n - f}
\leq (b - a)\,\norm{f_n - f}_\infty \longrightarrow 0 .
\]
وعامل الطول $(b - a)$ هو حيث يدخل تراص القطعة:
فعلى الفترات غير المتراصة يُنتج التقدير نفسه الحدَّ
عديم الفائدة $\infty\cdot0$، ويفشل الاستنتاج فعلًا
دون هيمنة --- فالنتوءات المسطحة $f_n =
\frac1n\mathbf 1_{\intcc{0}{n}}$ تتقارب بانتظام إلى $0$ على
$\intco{0}{\infty}$ ومع ذلك تُبقي $\int f_n = 1$ (انظر ملاحظة المزالق
أدناه)، وتفعل النتوءات المنزلقة في ملاحظة
\cref{ch:b2:integration} الشيء نفسه \omterm{def:b2:funcseq:def}{بتقارب نقطي}؛ \omterm{def:b2:funcseq:def}{فالتقارب
المنتظم} يتحكم في الارتفاعات، لا في العروض أبدًا.
\end{proof}

\begin{omfigure}
\begin{tikzpicture}
\begin{axis}[omaxis, width=8.6cm, height=5.4cm,
    xmin=0, xmax=1.05, ymin=0, ymax=1.9,
    xtick={0,1}, ytick={1}]
  \addplot[omDef!45, thick, domain=0:1.05, samples=120]
    {2*x*exp(-2*x^2)};
  \addplot[omDef!70, thick, domain=0:1.05, samples=120]
    {5*x*exp(-5*x^2)};
  \addplot[omDef, thick, domain=0:1.05, samples=160]
    {15*x*exp(-15*x^2)};
  \node[font=\small, omDef] at (axis cs:0.62,0.42) {$n=2$};
  \node[font=\small, omDef] at (axis cs:0.40,0.80) {$n=5$};
  \node[font=\small, omDef] at (axis cs:0.24,1.55) {$n=15$};
\end{axis}
\end{tikzpicture}

{\small النتوءات $g_n(x) = nx\,\eu^{-nx^2}$ في
\cref{exo:b2:funcseq:1}: تتقارب إلى $0$ عند كل نقطة،
لكن القمم (بارتفاع $\sim\sqrt{n/(2\eu)}$، منجرفةً نحو
$0$) تنمو بلا حدّ --- أي \omterm{def:b2:funcseq:def}{تقارب نقطي} مع
$\norm{g_n}_\infty \to \infty$، و$\int_0^1 g_n \to \frac12
\neq 0$: فالكتلة تختبئ تحت القمة المتحركة.}
\end{omfigure}

\begin{theorem}[الاشتقاق]\label{thm:b2:funcseq:differentiation}
لتكن $f_n$ من الصنف $C^1$ على فترة $I$، بحيث: تتقارب $(f_n')$
\emph{بانتظام} على $I$ (أو على كل قطعة من $I$) إلى دالة $g$ ما،
وتتقارب $(f_n(x_0))$ عند نقطة واحدة $x_0$. عندئذٍ تتقارب $(f_n)$
(بانتظام على القطع) إلى دالة $f$ من الصنف $C^1$، ويكون $f'
= g$: أي إن نهاية الدوال قابلة للاشتقاق حدًّا حدًّا.
\end{theorem}

\begin{proof}
نعرّف $f(x) = \lim f_n(x_0) + \int_{x_0}^x g$: وهذا مشروع، لأن $g$
متصلة --- فعلًا، $g$ هي النهاية \emph{المنتظمة} على
القطع للدوال المتصلة $f_n'$، ومن ثم تنطبق
\cref{thm:b2:funcseq:continuity}، وتكامل الدالة المتصلة
معرَّف جيدًا، ويكون بالمبرهنة الأساسية
للتفاضل والتكامل
\[
f'(x) = g(x) \qquad (x \in I) :
\]
فالنهاية المرشحة من الصنف $C^1$ وبالمشتق الصحيح
\emph{بحكم الإنشاء}، قبل برهان أي تقارب. وبالمبرهنة
الأساسية من جديد، $f_n(x) = f_n(x_0) + \int_{x_0}^x
f_n'$؛ وبالطرح،
\[
\abs{f_n(x) - f(x)} \leq \abs{f_n(x_0) - \lim f_n(x_0)}
+ \abs{x - x_0}\,\norm{f_n' - g}_{\infty} ,
\]
وهذا يؤول إلى $0$ بانتظام على كل قطعة. و$f$ من الصنف $C^1$
مع $f' = g$ بحكم الإنشاء.
\end{proof}

\begin{example}[لماذا تقع الفرضية على المشتقات]\label{ex:b2:funcseq:sqrtabs}
لتكن $F_n(x) = \sqrt{x^2 + \frac1n}$ على $\R$. وكل $F_n$ من الصنف
$C^1$ (بل $C^\infty$)، والتقارب إلى $\abs x$
منتظم على $\R$ كلها:
\[
0 \leq F_n(x) - \abs x
= \frac{(x^2 + \frac1n) - x^2}{\sqrt{x^2+\frac1n} + \abs x}
= \frac{1/n}{\sqrt{x^2 + \frac1n} + \abs x}
\leq \frac{1/n}{1/\sqrt n} = \frac{1}{\sqrt n} .
\]
ومع ذلك فالنهاية $\abs x$ غير قابلة للاشتقاق عند $0$: \omterm{def:b2:funcseq:def}{فالتقارب
المنتظم} \emph{للدوال}، مهما كان سريعًا، لا ينقل أي
قابلية للاشتقاق. والفشل ظاهر على المشتقات:
\[
F_n'(x) = \frac{x}{\sqrt{x^2 + \frac1n}}
\longrightarrow \begin{cases} 1 & x > 0,\\ 0 & x = 0,\\
-1 & x < 0, \end{cases}
\]
وهي نهاية نقطية غير متصلة، ومن ثم لا يمكن أن تتقارب $(F_n')$
بانتظام قرب $0$ (\cref{thm:b2:funcseq:continuity} من جديد).
والفكرة الختامية: تفترض
\cref{thm:b2:funcseq:differentiation} عمدًا
\omterm{def:b2:funcseq:def}{التقارب المنتظم} للمشتقات $f_n'$، لا للدوال $f_n$ ---
وهذا المثال هو السبب.
\end{example}

\section{متسلسلات الدوال}

\begin{definition}\label{def:b2:funcseq:series}
تتقارب متسلسلة دوال $\sum u_n$ \omterm{def:b2:funcseq:def}{نقطيًا} أو بانتظام حين تتقارب
مجاميعها الجزئية كذلك. وتتقارب
\emph{ناظميًا}\index{تقارب ناظمي} (على $X$) حين يكون $\sum
\norm{u_n}_\infty < \infty$. ويستلزم التقارب الناظمي التقاربَ
المنتظم (في \omterm{def:b2:nvs:banach}{فضاء باناخ} للدوال المحدودة:
\cref{thm:b2:nvs:absoluteconvergence})، الذي يستلزم النقطي؛
والاستلزامان تامّان.
\end{definition}

\begin{example}[متسلسلة واحدة، وثلاثة أحكام]\label{ex:b2:funcseq:threemodes}
لنأخذ $u_n(x) = \frac{x^n}{n}$ على $\intco{0}{1}$. \emph{\omterm{def:b2:funcseq:def}{نقطيًا}:}
تتقارب من أجل كل $x \in \intco01$ (بالمقارنة مع
المتسلسلة الهندسية). \emph{\omterm{def:b2:funcseq:series}{ناظميًا} على $\intcc{0}{a}$، مع $a < 1$:}
$\norm{u_n}_{\infty,\intcc0a} = \frac{a^n}{n}$، وهي \omterm{def:b2:series:summable}{قابلة للجمع}.
\emph{وليست ناظمية على $\intco{0}{1}$:}
$\norm{u_n}_{\infty,\intco01} = \frac1n$، و$\sum\frac1n$
تتباعد. \emph{وليست حتى منتظمة على $\intco{0}{1}$:} إذ
يقاوم الباقي قرب $1$،
\[
R_N(x) = \sum_{n>N}\frac{x^n}{n}
\geq \sum_{n=N+1}^{2N}\frac{x^n}{n}
\geq \frac{N\,x^{2N}}{2N} = \frac{x^{2N}}{2}
\xrightarrow[x\to1^-]{} \frac12 ,
\]
ومن ثم $\sup_{\intco01}\abs{R_N} \geq \frac12$ من أجل كل $N$.
والفكرة الختامية: تتعايش الأحكام الأربعة بسلام --- فالمجموع
$-\ln(1-x)$ \omterm{def:b2:metric:continuity}{متصل} على $\intco{0}{1}$ لأن
الاتصال لا يحتاج إلا الانتظام \emph{قرب كل نقطة}، أي
على القطع $\intcc0a$؛ والانفجار عند الحافة حقٌّ
للمجموع.
\end{example}

\begin{theorem}[النقل من أجل المتسلسلات]\label{thm:b2:funcseq:seriestransfer}
إذا تقاربت $\sum u_n$ بانتظام (مثلًا \omterm{def:b2:funcseq:series}{ناظميًا}) على المجموعة
المعنية: فينتقل اتصال كل $u_n$ عند $a$ إلى المجموع؛ ويمكن المكاملة
على قطعة حدًّا حدًّا؛ وإذا تقاربت $\sum u_n(x_0)$
وتقاربت $\sum u_n'$ بانتظام على القطع، فإن
المجموع من الصنف $C^1$ ومشتقه $\sum u_n'$.
\end{theorem}

\begin{proof}
كل شيء هو المبرهنة الموافقة مطبَّقةً على المجاميع
الجزئية $S_N = \sum_{n\leq N}u_n$، وهي مجاميع منتهية لدوال
تتمتع بالانتظام المعني. الاتصال: كل $S_N$
متصلة عند $a$ و$S_N \to \sum u_n$ بانتظام:
\cref{thm:b2:funcseq:continuity}. المكاملة: على القطعة،
\[
\int_a^b \sum_{n\geq0} u_n
= \lim_N \int_a^b S_N
= \lim_N \sum_{n=0}^{N}\int_a^b u_n
= \sum_{n\geq0}\int_a^b u_n
\]
حسب \cref{thm:b2:funcseq:integration} (المساواة الأولى) و
بخطية التكامل (الثانية). الاشتقاق: الدوال
$S_N$ من الصنف $C^1$، وتتقارب $S_N(x_0)$، وتتقارب $S_N' =
\sum_{n\leq N}u_n'$ بانتظام على القطع:
فتعطي \cref{thm:b2:funcseq:differentiation} أن المجموع
من الصنف $C^1$ بالمشتق $\lim S_N' = \sum u_n'$.
\end{proof}

\begin{remark}[مزالق شائعة]
أربعة فخاخ، وقد رُئيت كلها في أوراق الامتحان. \emph{(1)
النهايات العليا المتحقَّق منها نصفًا:} تقييم $f_n$ على امتداد متتالية
مختارة جيدًا $x_n$ لا يحدّ $\norm{f_n - f}_\infty$ إلا
من \emph{أسفل} --- وهذا يكفي لدحض الانتظام (كما في
\cref{ex:b2:funcseq:xn})، ولا يكفي أبدًا لبرهانه؛ ولبرهانه، حُدّ
النهاية العليا بحساب صالح من أجل \emph{كل} $x$. \emph{(2)
الانتظام على المجموعة الخطأ:} كثيرًا ما يصحّ \omterm{def:b2:funcseq:series}{التقارب الناظمي} أو المنتظم
على كل $\intcc{-a}{a}$ أو $\intco\delta\infty$
لكنه يفشل على الاتحاد المفتوح؛ وليس ذلك عائقًا --- فالاتصال
وقابلية الاشتقاق محليان، ومن ثم يعطيهما انضباط قطعةً قطعة
في \cref{ex:b2:funcseq:zeta} على المجموعة \omterm{def:b2:metric:topology}{المفتوحة} كلها. \emph{(3)
المكاملة على غير القطع:}
\cref{thm:b2:funcseq:integration} قول عن
القطع؛ فعلى $\intco0\infty$، لا يمنع \omterm{def:b2:funcseq:def}{التقارب المنتظم}
هروب الكتلة إلى اللانهاية (فالدوال $f_n =
\frac1n\mathbf 1_{\intcc{0}{n}}$ تتقارب بانتظام إلى $0$،
مع $\int f_n = 1$) --- فاستعمل التقارب المهيمَن عليه هناك.
\emph{(4) اشتقاق النهاية:}
\cref{ex:b2:funcseq:sqrtabs}؛ فالفرضية على المشتقات $(f_n')$، ولا يمكن لأي معدل تقارب للدوال $(f_n)$ أن يعوّضها.
\end{remark}

\begin{example}[دالة ريمان $\zeta$]\label{ex:b2:funcseq:zeta}
تتقارب $\zeta(s) = \sum_{n\geq1} n^{-s}$ \omterm{def:b2:funcseq:series}{ناظميًا} على كل نصف
مستقيم $\intco{a}{+\infty}$، مع $a > 1$ (لأن $\norm{n^{-s}}_\infty =
n^{-a}$، وهي \omterm{def:b2:series:summable}{قابلة للجمع}): ومن ثم فإن $\zeta$ متصلة على $\intoo{1}{+\infty}$؛
وبالاشتقاق حدًّا حدًّا (فالمتسلسلة المشتقة $\sum -\ln n\; n^{-s}$
تتقارب \omterm{def:b2:funcseq:series}{ناظميًا} هي أيضًا على $\intco{a}{\infty}$)، تكون $\zeta$ من الصنف $C^1$
--- وبالتكرار، من الصنف $C^\infty$ --- مع $\zeta'(s) = -\sum \frac{\ln
n}{n^s}$. ولاحظ الانضباط: يُتحقَّق من \omterm{def:b2:funcseq:series}{التقارب الناظمي} على
\emph{أنصاف} المستقيمات الجزئية، لا على المفتوح $\intoo{1}{\infty}$
نفسه، حيث يفشل.
\end{example}

\begin{example}[متسلسلة لوغاريتمية، منفَّذة إلى النهاية]\label{ex:b2:funcseq:logseries}
لتكن $F(x) = \sum_{n\geq1} \frac{\eu^{-nx}}{n}$ على
$\intoo{0}{\infty}$. وكل حدّ محدود على $\intco{\delta}
\infty$ بالمقدار $\frac{\eu^{-n\delta}}{n} \leq \eu^{-n\delta}$، وهو حدّ
متسلسلة هندسية متقاربة: \omterm{def:b2:funcseq:series}{فالتقارب ناظمي} على كل
$\intco\delta\infty$، ومن ثم فإن $F$ متصلة على
$\intoo{0}{\infty}$. والمتسلسلة المشتقة $\sum -\eu^{-nx}$
متقاربة \omterm{def:b2:funcseq:series}{ناظميًا} كذلك على $\intco\delta\infty$
($\norm{\eu^{-nx}}_{\infty,\intco\delta\infty} =
\eu^{-n\delta}$)، ومن ثم فإن $F$ من الصنف $C^1$ بمشتق هندسي:
\[
F'(x) = -\sum_{n\geq1}\eu^{-nx}
= \frac{-\eu^{-x}}{1 - \eu^{-x}}
= \frac{-1}{\eu^{x} - 1} .
\]
وبالتكرار، تكون $F$ من الصنف $C^\infty$. وبمكاملة $F'$ (فكلٌّ من $F$ و$x
\mapsto -\ln(1 - \eu^{-x})$ ينعدم عند $+\infty$ ولهما
المشتق نفسه على $\intoo0\infty$):
\[
F(x) = -\ln\bigl(1 - \eu^{-x}\bigr),
\]
وهي المتسلسلة اللوغاريتمية عند $t = \eu^{-x}$. والفكرة الختامية: حين $x
\to 0^+$، $F(x) = -\ln(x + O(x^2)) = \ln\frac1x + O(x)$ --- فتتباعد
المتسلسلة لوغاريتميًا عند الحدّ، تمامًا مثل
المتسلسلة التوافقية التي تصيرها عند $x = 0$؛ \omterm{def:b2:funcseq:series}{والتقارب الناظمي} على
$\intco\delta\infty$ دون $\intoo0\infty$ هو
العَرَض.
\end{example}

\begin{method}[كيف نبرهن على التقارب المنتظم أو ندحضه]\label{met:b2:funcseq:uniform}
من أجل $f_n \to f$ \omterm{def:b2:funcseq:def}{نقطيًا} على $X$:
\begin{enumerate}
  \item احسب أو حُدّ $\norm{f_n - f}_\infty$: ادرس
        الدالة $x \mapsto \abs{f_n(x) - f(x)}$ (بالمشتق و
        الرتابة) لتحديد موضع أعظميتها؛ فحدٌّ صالح من أجل
        كل $x$ ويؤول إلى $0$ يبرهن على الانتظام.
  \item \emph{وللدحض}: أبرِز نقاطًا $x_n$ تحقق
        $\abs{f_n(x_n) - f(x_n)} \not\to 0$ (وكثيرًا ما يتتبع $x_n$
        النتوء المتحرك، كما في \cref{exo:b2:funcseq:1})؛ أو
        استشهد بمبرهنة نقل بالعكس المنطقي --- نهاية غير
        متصلة لدوال متصلة
        (\cref{ex:b2:funcseq:threshold})، أو $\int f_n
        \not\to \int f$ على قطعة.
  \item ومن أجل المتسلسلات، جرّب \omterm{def:b2:funcseq:series}{التقارب الناظمي} أولًا
        ($\sum\sup\abs{u_n} < \infty$)؛ فإذا فشل شموليًا،
        فاختبره على القطع الجزئية المهمة
        (\cref{ex:b2:funcseq:threemodes})؛ وإذا فشل
        في كل مكان، فقد يصحّ \omterm{def:b2:funcseq:def}{التقارب المنتظم} مع ذلك عبر
        حدّ باقي المتناوبات
        (\cref{exo:b2:funcseq:4}) أو عبر الجمع بالتجزئة.
\end{enumerate}
\end{method}

\section{مبرهنة فايرشتراس في التقريب}

\begin{theorem}[فايرشتراس، عبر برنشتاين]\label{thm:b2:funcseq:weierstrass}
كل دالة متصلة $f \colon \intcc{0}{1} \to \R$ نهايةٌ منتظمة
لكثيرات حدود --- وبالتحديد، \emph{لكثيرات حدود
برنشتاين}\index{كثيرات حدود برنشتاين}\index{مبرهنة فايرشتراس
في التقريب} الخاصة بها
\[
B_n(f)(x) = \sum_{k=0}^{n} f\Bigl(\frac kn\Bigr)\binom nk x^k
(1-x)^{n-k} .
\]
\end{theorem}

\begin{proof}
نثبّت $x \in \intcc{0}{1}$ ونضع $p_k(x) = \binom nk x^k(1 -
x)^{n-k}$. وثلاث متطابقات ثنائية، نحصل عليها بتقييم
$(x + y)^n$ ومشتقيه بالنسبة إلى $x$ عند $y = 1 - x$:
\[
\sum_k p_k = 1,
\qquad
\sum_k k\,p_k = nx,
\qquad
\sum_k k(k-1) p_k = n(n-1)x^2 .
\]
وبالتفصيل: $(x+y)^n = \sum_k\binom nk x^ky^{n-k}$ عند $y = 1-x$
هي الأولى؛ وبالاشتقاق بالنسبة إلى $x$،
\[
n(x+y)^{n-1} = \sum_k k\binom nk x^{k-1}y^{n-k} ,
\]
ثم بالضرب في $x$ ووضع $y = 1 - x$ نحصل على
الثانية؛ وبالاشتقاق مرتين والضرب في $x^2$ نحصل
على الثالثة. وبنشر $(k - nx)^2 = k(k-1) + k(1 - 2nx) +
n^2x^2$ وجمع الثلاث:
\[
\sum_k (k - nx)^2 p_k
= n(n-1)x^2 + nx(1 - 2nx) + n^2x^2
= nx(1 - x) \leq \frac n4 ,
\]
وهي \emph{متطابقة التباين}.

ونقدّر الآن، مستعملين $\sum p_k = 1$:
\[
\abs{B_n(f)(x) - f(x)}
\leq \sum_{k} \Bigl| f\Bigl(\frac kn\Bigr) - f(x)\Bigr|\, p_k(x)
= \Sigma_{\text{قريب}} + \Sigma_{\text{بعيد}} ,
\]
بالفصل حسب $\abs{\frac kn - x} \leq \delta$ أو لا.
ومن أجل $\varepsilon > 0$ معطى، يعطي الاتصال المنتظم للدالة $f$ (هاينه)
عددًا $\delta$ يحقق $\Sigma_{\text{قريب}} \leq \varepsilon$. وأما المجموع
البعيد، مع $M = \norm f_\infty$: فبمتطابقة التباين وبحيلة
العدّ عند تشيبيشيف،
\[
\Sigma_{\text{بعيد}} \leq 2M \sum_{\abs{k - nx} > n\delta} p_k
\leq 2M\,\frac{\sum_k (k - nx)^2 p_k}{n^2\delta^2}
\leq \frac{2M}{4 n \delta^2} \cdot 1
= \frac{M}{2n\delta^2}
\xrightarrow[n\to\infty]{} 0 ,
\]
بانتظام في $x$. ومن ثم $\norm{B_n(f) - f}_\infty \leq \varepsilon +
\frac{M}{2n\delta^2} \leq 2\varepsilon$ من أجل $n$ الكبيرة.
\end{proof}

\begin{remark}
بالتعويض الأفيني تصحّ المبرهنة على أي قطعة
$\intcc{a}{b}$. وهي تفشل على $\R$ (فالنهاية المنتظمة لكثيرات حدود على
$\R$ كثير حدود: \cref{exo:b2:funcseq:8}). أما القراءة الاحتمالية
--- فالمقدار $B_n(f)(x)$ هو القيمة المتوقعة للدالة $f$ عند متوسط
ثنائي، وحدّ التباين هو متراجحة تشيبيشيف --- فتُصاغ بأمانة في \cref{ch:b2:genfun}.
\end{remark}

\begin{omfigure}
\begin{tikzpicture}
\begin{axis}[omaxis, width=8.6cm, height=6.0cm,
    xmin=0, xmax=1.05, ymin=0, ymax=1.08,
    xtick={0,1}, ytick={1}]
  \addplot[omThm, very thick, domain=0:1, samples=80] {x^2};
  \addplot[omDef!45, thick, domain=0:1, samples=2] {x};
  \addplot[omDef!70, thick, domain=0:1, samples=80]
    {x^2 + x*(1-x)/2};
  \addplot[omDef, thick, domain=0:1, samples=80]
    {x^2 + x*(1-x)/6};
  \node[font=\small, omDef] at (axis cs:0.30,0.37) {$B_1f$};
  \node[font=\small, omDef] at (axis cs:0.52,0.46) {$B_2f$};
  \node[font=\small, omDef] at (axis cs:0.76,0.56) {$B_6f$};
  \node[font=\small, omThm] at (axis cs:0.88,0.68) {$f$};
\end{axis}
\end{tikzpicture}

{\small تقريب برنشتاين للدالة $f(x) = x^2$ (بالأحمر)، باستعمال
الصيغة الدقيقة $B_nf = x^2 + \frac{x(1-x)}{n}$ في
\cref{exo:b2:funcseq:7}: $B_1f$ هو الوتر، وكل مضاعفة
للعدد $n$ تنصّف الفجوة. موثوق لكنه بطيء --- وهو التشبّع من الرتبة
$\frac1n$ الذي تجعله مبرهنة فورونوفسكايا (مسألة نهاية الأسبوع)
دقيقًا.}
\end{omfigure}

\begin{remark}[أين يُستعمَل هذا]
مبرهنة فايرشتراس في التقريب هي مبرهنة الكثافة في التحليل
الكلاسيكي: فهي تجعل $C(\intcc ab)$ قابلًا للفصل، وتتيح التحقق من
المتطابقات التكاملية على كثيرات الحدود وحدها (مسائل العزوم)، و
تقوم تحت النسخة المثلثية المبرهَن عليها في فصل فورييه
عبر نواة فييير. وتستخرج مسألة نهاية الأسبوع في هذا الفصل
المضمونَ الكمّي لبرهان برنشتاين
--- أي معدلات التقارب التي يحكمها مقياس الاتصال ---
ثم تعزل ما جعله يعمل حقًا، في مبرهنة كوروفكين:
الإيجابية مع ثلاث دوال اختبار. ويعمّم مجلد السنة
الثالثة قول الكثافة على الجبور الجزئية الكيفية
(ستون--فايرشتراس) وعلى الفضاءات المتراصة.
\end{remark}

\begin{example}[تقريب مضلّعي، بمعدل]\label{ex:b2:funcseq:polygonal}
من أجل دالة $f$ ليبشيتزية بالثابت $L$ على $\intcc{0}{1}$، لتكن $I_nf$
الدالة المستكمِلة الأفينية على قطع عند العقد $\frac kn$. وعلى
خلية $\intcc{\frac kn}{\frac{k+1}n}$، يقع كلٌّ من $f(x)$ و
$I_nf(x)$ بين القيمتين الحدّيتين اللتين يمكن لدالة ليبشيتزية بالثابت $L$
أن تأخذهما بمعلومية القيمتين العقديتين، ومن ثم من أجل $x$ في
الخلية، وبكتابة $x_k = \frac kn$:
\[
\abs{I_nf(x) - f(x)}
\leq \abs{I_nf(x) - f(x_k)} + \abs{f(x_k) - f(x)}
\leq L\,\abs{x - x_k} + L\,\abs{x - x_k}
\leq \frac{2L}{n}
\]
(فالدالة المستكمِلة نفسها ليبشيتزية بالثابت $L$ على الخلية: إذ إن
ميلها نسبةُ فروق للدالة $f$). ومنه $\norm{I_nf -
f}_\infty \leq \frac{2L}{n}$: فالتقريب المضلّعي للدوال
الليبشيتزية يتقارب بسرعة $\frac1n$ ---
\emph{أسرع} من معدل برنشتاين $\frac{1}{\sqrt n}$ من أجل
الصنف نفسه (مسألة نهاية الأسبوع، الجزء الثاني). والفكرة الختامية:
المضلّع يستكمل لكنه ليس أملس، وبرنشتاين أملس
لكنه بطيء؛ فلا غداء مجاني بين انتظام
المقرِّب والسرعة --- وهي مقايضة تجعلها نتائج التشبّع
في مسألة نهاية الأسبوع دقيقة.
\end{example}

\begin{remark}[آفاق داخل هذا المجلد]
\omterm{def:b2:funcseq:def}{التقارب المنتظم} حصان عمل هذا الكتاب من الآن فصاعدًا. فيجري
فصل متسلسلات القوى كله على \omterm{def:b2:funcseq:series}{التقارب الناظمي} على
الأقراص الجزئية المتراصة --- وكل مبرهنة حدًّا حدًّا هناك
حالةٌ خاصة من مبرهنات النقل في هذا الفصل. ويعيش فصل
فورييه طابقًا أعلى: فمجاميعه الجزئية $S_N$ تفشل
بالضبط حيث يحذّر هذا الفصل من أنها قد تفشل (\omterm{def:b2:funcseq:def}{نقطيًا} لا
بانتظام عند القفزات)، وتنجح متوسطات فييير فيه بميكانيك
$3\varepsilon$ نفسه الذي برهن على
\cref{thm:b2:funcseq:continuity}. ويعرّف فصل المعادلات
التفاضلية $\eu^{tA}$ بمتسلسلة متقاربة \omterm{def:b2:funcseq:series}{ناظميًا} و
يشتقها حدًّا حدًّا --- وهي حرفيًا
\cref{thm:b2:funcseq:seriestransfer} مطبَّقة على عناصر
المصفوفات. وحين يخالجك الشك لاحقًا في الكتاب عن ``لماذا يجوز لنا
أن نفعل هذا''، فالجواب عادةً مبرهنة من هذا الفصل.
\end{remark}

\section{تمارين}

\begin{exercise}[$\star$]\label{exo:b2:funcseq:1}
ادرس \omterm{def:b2:funcseq:def}{التقارب النقطي} والمنتظم على $\intcc{0}{1}$، ثم
على $\intcc{0}{a}$ (مع $a < 1$) أو على $\intco{\delta}{1}$ حسب الاقتضاء، من أجل:
\[
f_n(x) = \frac{x}{1 + nx},
\qquad
g_n(x) = n x\,\eu^{-n x^2},
\qquad
h_n(x) = x^n(1 - x^n).
\]
\end{exercise}

\begin{exercise}[$\star$]\label{exo:b2:funcseq:2}
برهن على أن $\displaystyle\int_0^1 g_n \not\to \int_0^1 \lim g_n$
من أجل $g_n$ في \cref{exo:b2:funcseq:1}، ووفّق ذلك مع
\cref{thm:b2:funcseq:integration}.
\end{exercise}

\begin{exercise}[$\star$]\label{exo:b2:funcseq:3}
برهن على أن $S(x) = \sum_{n\geq1} \dfrac{x^n}{n^2}$ متصلة على
$\intcc{-1}{1}$، وأن $S$ من الصنف $C^1$ على $\intoo{-1}{1}$ مع
$S'(x) = -\frac{\ln(1-x)}{x}$ من أجل $0 < \abs x < 1$.
\end{exercise}

\begin{exercise}[$\star\star$]\label{exo:b2:funcseq:4}
لتكن $F(x) = \sum_{n \geq 0} \dfrac{(-1)^n}{n + x}$ على
$\intoo{0}{+\infty}$. برهن على \omterm{def:b2:funcseq:def}{التقارب المنتظم} (لا الناظمي) على
$\intco{\delta}{\infty}$ عبر حدّ باقي المتسلسلات
المتناوبة، وعلى الاتصال، وعلى المعادلة الوظيفية $F(x) + F(x + 1) =
\frac1x$.
\end{exercise}

\begin{exercise}[$\star\star$]\label{exo:b2:funcseq:5}
(ديني) لتكن $f_n \colon K \to \R$ متصلة على \omterm{def:b2:metric:def}{فضاء متري}
\emph{متراص}، مع $f_n \to f$ \emph{\omterm{def:b2:funcseq:def}{نقطيًا}}، و$f$ متصلة،
و$(f_n(x))$ \emph{متناقصة} في $n$ من أجل كل $x$. برهن على أن
\omterm{def:b2:funcseq:def}{التقارب منتظم}. \emph{(من أجل $\varepsilon$ معطى، تتزايد المجموعات
\omterm{def:b2:metric:topology}{المفتوحة} $U_n = \{x : f_n(x) - f(x) < \varepsilon\}$ وتغطي
$K$؛ فاستخرج تغطية جزئية منتهية ---
\cref{thm:b2:metric:borellebesgue}.)}
\end{exercise}

\begin{exercise}[$\star\star$]\label{exo:b2:funcseq:6}
برهن على أن $\displaystyle\lim_{n\to\infty} \int_0^1
\frac{n\,f(x)}{1 + n^2x^2}\,\dd x = \frac{\pi}{2} f(0)$ من أجل كل
دالة متصلة $f$ على $\intcc{0}{1}$. \emph{(عوّض $u = nx$؛ وافصل
$f(0)$؛ وهيمن.)}
\end{exercise}

\begin{exercise}[$\star\star$]\label{exo:b2:funcseq:7}
احسب \omterm{thm:b2:funcseq:weierstrass}{كثيرات حدود برنشتاين} للدالة $f(x) = x^2$ صراحةً و
تحقق من الخطأ المنتظم $\norm{B_n f - f}_\infty =
O\bigl(\frac1n\bigr)$ المتنبَّأ به في برهان
\cref{thm:b2:funcseq:weierstrass} --- وهو هنا $\frac{x(1 -
x)}{n}$ بالضبط عند كل نقطة.
\end{exercise}

\begin{exercise}[$\star\star$]\label{exo:b2:funcseq:8}
برهن على أنه إذا تقاربت كثيرات حدود $P_n$ بانتظام \emph{على كل
$\R$} إلى $f$، فإن $f$ كثير حدود. \emph{(من أجل $m, n$ الكبيرة،
يكون $P_n - P_m$ كثير حدود محدودًا على $\R$، ومن ثم ثابتًا؛ فتستقر
المتتالية إلى حدّ ثوابت.)}
\end{exercise}

\begin{exercise}[$\star\star\star$]\label{exo:b2:funcseq:9}
(دالة متصلة غير قابلة للاشتقاق في أي نقطة --- موجَّه) لتكن
$\varphi$ المسافة إلى أقرب عدد صحيح (وهي دورية بالدور $1$، و
$\norm{\varphi}_\infty = \frac12$، وليبشيتزية بالثابت $1$) ولتكن
\[
W(x) = \sum_{n=0}^{\infty} \Bigl(\frac{3}{4}\Bigr)^{\!n}
\varphi(4^n x) .
\]
برهن على: (أ) أن $W$ متصلة على $\R$ (\omterm{def:b2:funcseq:series}{بالتقارب الناظمي})؛ (ب) أنه من أجل
كل $x$ وكل $m$، وباختيار $h_m = \pm\frac12\cdot 4^{-m}$
بالإشارة التي تجعل $\varphi$ أفينية على القطعة من $4^m x$
إلى $4^m(x + h_m)$، تحقق نسبة الفروق
\[
\Bigl|\frac{W(x + h_m) - W(x)}{h_m}\Bigr| \geq 3^m -
\sum_{n<m} 3^n \geq \frac{3^m + 1}{2} \xrightarrow[m\to\infty]{}
\infty
\]
(فالحدود $n > m$ تنعدم بالدورية؛ والحدّ $n = m$ يسهم
بالمقدار $3^m$ بالضبط؛ والحدود $n < m$ محدودة بخاصية ليبشيتز).
واختم بأن $W$ غير قابلة للاشتقاق في أي نقطة.
\end{exercise}

\begin{exercise}[$\star$]\label{exo:b2:funcseq:10}
لتكن $u_n(x) = (-1)^n x^n(1 - x)$ على $\intcc{0}{1}$. بيّن أن
$\sum u_n$ تتقارب \omterm{def:b2:funcseq:def}{نقطيًا} على $\intcc{0}{1}$ واحسب
مجموعها؛ وبيّن أن \omterm{def:b2:funcseq:def}{التقارب منتظم} على $\intcc{0}{1}$
\emph{(حُدّ الباقي $R_N(x) = \sum_{n > N} u_n(x)$، وهو ذيل
هندسي، بحدّه الأول وعظّم $x^{N+1}(1-x)$)}
لكنه \emph{ليس} \omterm{def:b2:funcseq:series}{ناظميًا} \emph{(احسب $\norm{u_n}_\infty$)}:
\omterm{def:b2:funcseq:def}{فالتقارب المنتظم} أضعف تمامًا من \omterm{def:b2:funcseq:series}{التقارب الناظمي}.
وقارِن بالمتسلسلة $\sum x^n(1-x)$، التي مجموعها غير \omterm{def:b2:metric:continuity}{متصل} عند
$1$: فهناك يفشل حتى الانتظام.
\end{exercise}

\begin{exercise}[$\star\star$]\label{exo:b2:funcseq:11}
لتكن $f_n \to f$ بانتظام على \omterm{def:b2:metric:def}{فضاء متري} $X$، وكل $f_n$
متصلة، وليكن $x_n \to x$ في $X$. برهن على أن $f_n(x_n) \to
f(x)$. وبيّن بمثال على $X = \intcc{0}{1}$ أن \omterm{def:b2:funcseq:def}{التقارب النقطي}
لا يكفي، حتى مع $f$ متصلة
\emph{(استعمل النتوءات $g_n$ في \cref{exo:b2:funcseq:1} و$x_n =
\frac{1}{\sqrt{2n}}$)}.
\end{exercise}

\begin{exercise}[$\star\star\star$]\label{exo:b2:funcseq:12}
(معادلة تكاملية لفولتيرا بالمتسلسلات) من أجل $f \in
C(\intcc{0}{1})$ نعرّف $Tf(x) = \int_0^x f(t)\,\dd t$.
\begin{enumerate}
  \item بيّن بالتراجع أنه من أجل $n \geq 1$:
        \[
        T^n f(x) = \int_0^x \frac{(x -
        t)^{n-1}}{(n-1)!}\,f(t)\,\dd t,
        \qquad
        \norm{T^n f}_\infty \leq \frac{\norm f_\infty}{n!} .
        \]
  \item استنتج أن $S = \sum_{n\geq0} T^n f$ تتقارب \omterm{def:b2:funcseq:series}{ناظميًا}
        على $\intcc{0}{1}$ وتحلّ المعادلة التكاملية $S = f
        + TS$.
  \item تحقق من أن $S(x) = f(x) + \int_0^x \eu^{x-t}f(t)\,\dd
        t$ يحلّ المعادلة نفسها، وبرهن على وحدانية
        الحلول المتصلة \emph{(إذا كان $S = TS$ فإن
        $\norm{S}_\infty \leq \norm{T^nS}_\infty \to 0$)}:
        واختم بالصورة المغلقة للمجموع.
\end{enumerate}
\end{exercise}

\section{مسألة: معدلات التقريب ومبرهنة كوروفكين}

\begin{problem}\label{pb:b2:funcseq:1}
يخفي برهان برنشتاين للمبرهنة \cref{thm:b2:funcseq:weierstrass} كنزين.
أولًا، إنه \emph{كمّي}: فسرعة $B_nf
\to f$ يحكمها مقياس اتصال $f$، بالمعدل
الأمثل المبلوغ عند $\abs{x - \frac12}$. وثانيًا، إنه
\emph{بنيوي}: فكل ما أهمّ هو أن $B_n$ مؤثر خطي
موجب يسلك سلوكًا حسنًا على $1$ و$x$ و$x^2$ ---
وهذه الملاحظة، معزولةً، هي \emph{مبرهنة
كوروفكين}\index{مبرهنة كوروفكين}. وتبرهن هذه المسألة على الاثنين،
وتُختم بالمقارب الدقيق عند فورونوفسكايا. وفي كل ما يلي،
$f \in C(\intcc{0}{1})$، $M = \norm f_\infty$، $p_k(x) = \binom
nk x^k(1-x)^{n-k}$، ويرمز $e_j$ إلى $x \mapsto x^j$.

\medskip
\noindent\textbf{الجزء الأول --- مؤثر برنشتاين.}
\begin{enumerate}
  \item بيّن أن $B_n$ خطي و\emph{موجب} ($f \geq 0
        \Rightarrow B_nf \geq 0$)، ومن ثم رتيب ($f \leq g
        \Rightarrow B_nf \leq B_ng$)، مع
        $\norm{B_nf}_\infty \leq \norm f_\infty$، وأن
        $B_nf$ يستكمل $f$ عند الطرفين معًا.
  \item أعد استنباط المتطابقات $B_n e_0 = e_0$ و$B_n e_1 =
        e_1$ و$B_n e_2 = e_2 + \frac{e_1 - e_2}{n}$
        \emph{(اشتق $(x + y)^n$ مرتين وضع $y = 1 -
        x$)}.
  \item استنتج متطابقة التباين $\sum_k \bigl(\frac kn -
        x\bigr)^2 p_k(x) = \frac{x(1-x)}{n}$ ثم، بكوشي--شوارتز،
        حدَّ العزم الأول
        \[
        \sum_{k=0}^{n}\Bigl|\frac kn - x\Bigr|\,p_k(x)
        \leq \sqrt{\frac{x(1-x)}{n}} \leq \frac{1}{2\sqrt n} .
        \]
  \item بيّن أنه إذا كانت $f$ محدبة فإن $B_nf \geq f$ على
        $\intcc{0}{1}$ \emph{(بمتراجحة جنسن المنتهية من أجل
        الأوزان $p_k(x)$)}.
  \item (حدّ العدّ عند تشيبيشيف، معادًا صياغته) من أجل $\delta > 0$
        بيّن أن
        \[
        \sum_{\abs{k/n - x} > \delta} p_k(x)
        \leq \frac{x(1-x)}{n\delta^2}
        \leq \frac{1}{4n\delta^2} ,
        \]
        وأعطِ القراءة الاحتمالية: $B_nf(x)$ يأخذ متوسط
        $f$ على متوسط عينة ثنائية يتركّز عند $x$.
\end{enumerate}

\medskip
\noindent\textbf{الجزء الثاني --- المعدلات: مقياس الاتصال.}
من أجل $\delta > 0$ نضع $\omega(\delta) = \sup\{\abs{f(s) - f(t)}
: s, t \in \intcc{0}{1},\ \abs{s - t} \leq \delta\}$.
\begin{enumerate}[resume]
  \item بيّن أن: $\omega$ منته وغير متناقص، وأن
        $\omega(\delta) \to 0$ حين $\delta \to 0^+$ (هاينه)، وأنه
        تحت جمعي ($\omega(\delta_1 + \delta_2) \leq
        \omega(\delta_1) + \omega(\delta_2)$)، وأن
        $\omega(\lambda\delta) \leq (1 +
        \lambda)\,\omega(\delta)$ من أجل كل $\lambda > 0$.
  \item برهن على التقدير الرئيسي، من أجل كل $\delta > 0$:
        \[
        \abs{B_nf(x) - f(x)}
        \leq \sum_k \omega\Bigl(\Bigl|\frac kn -
        x\Bigr|\Bigr)p_k(x)
        \leq \Bigl(1 + \frac1\delta\sum_k\Bigl|\frac kn -
        x\Bigr|p_k(x)\Bigr)\,\omega(\delta) .
        \]
  \item اختر $\delta = n^{-1/2}$ واختم \emph{مبرهنة
        فايرشتراس الكمّية}:
        \[
        \norm{B_nf - f}_\infty \leq
        \frac32\,\omega\Bigl(\frac{1}{\sqrt n}\Bigr)
        \xrightarrow[n\to\infty]{} 0 .
        \]
  \item استنتج المعدلات: $\norm{B_nf - f}_\infty \leq
        \frac{3L}{2\sqrt n}$ من أجل $f$ ليبشيتزية بالثابت $L$، و$\leq
        \frac32 C n^{-\alpha/2}$ من أجل $f$ هولدرية بالأس $\alpha$
        ($\abs{f(s) - f(t)} \leq C\abs{s-t}^\alpha$).
  \item (المثال الأمثل --- متطابقة ثنائية) من أجل $m
        \geq 1$ برهن على أن
        \[
        \sum_{k=m+1}^{2m} (k - m)\binom{2m}{k}
        = \frac{m}{2}\binom{2m}{m},
        \qquad\text{ومنه}\qquad
        \sum_{k=0}^{2m}\abs{k - m}\binom{2m}{k}
        = m\binom{2m}{m}
        \]
        \emph{(استعمل $k\binom{2m}k = 2m\binom{2m-1}{k-1}$ و
        تناظر السطر الثنائي، الذي يعطي
        $\sum_{j=m}^{2m-1}\binom{2m-1}{j} = 2^{2m-2}$)}.
  \item من أجل $f(t) = \abs{t - \frac12}$ استنتج القيمة الدقيقة
        ومقارباتها (بالمعامل الثنائي المركزي،
        \cref{ex:b2:comparison:centralbinomial}):
        \[
        B_{2m}f\Bigl(\frac12\Bigr) - f\Bigl(\frac12\Bigr)
        = \frac{\binom{2m}{m}}{2\cdot4^{m}}
        \;\sim\; \frac{1}{2\sqrt{\pi m}} :
        \]
        فالمعدل $\omega(n^{-1/2})$ في السؤال 8 مبلوغ
        (إلى حدّ ثابت) --- ومن أجل $f$ المتصلة فحسب،
        يكون معدل برنشتاين $n^{-1/2}$ أمينًا.
\end{enumerate}

\medskip
\noindent\textbf{الجزء الثالث --- مبرهنة كوروفكين.} لتكن $(L_n)$
متتالية من المؤثرات \emph{الخطية الموجبة} من
$C(\intcc{0}{1})$ إلى نفسه بحيث $L_ne_j \to e_j$
بانتظام من أجل $j = 0, 1, 2$.
\begin{enumerate}[resume]
  \item بيّن أن المؤثر الخطي الموجب $L$ رتيب و
        يحقق $\abs{Lf} \leq L\abs f$ \omterm{def:b2:funcseq:def}{نقطيًا}.
  \item بيّن أنه من أجل كل $\varepsilon > 0$ يوجد $\delta > 0$
        بحيث من أجل \emph{كل} $s, x \in \intcc{0}{1}$:
        \[
        \abs{f(s) - f(x)} \leq \varepsilon +
        \frac{2M}{\delta^2}(s - x)^2
        \]
        \emph{(عالج $\abs{s - x} \leq \delta$ بهاينه و
        $\abs{s-x} > \delta$ بالحدّ الخام $2M$)}.
  \item ثبّت $x$، وطبّق $L_n$ على متراجحة السؤال 13
        في المتغير $s$، واستنبط
        \[
        \abs{L_nf(x) - f(x)\,L_ne_0(x)}
        \leq \varepsilon\,L_ne_0(x) + \frac{2M}{\delta^2}
        \bigl(L_ne_2(x) - 2x\,L_ne_1(x) + x^2 L_ne_0(x)\bigr).
        \]
  \item بيّن أن $\sup_x \bigl(L_ne_2(x) - 2x\,L_ne_1(x) + x^2
        L_ne_0(x)\bigr) \to 0$، ثم ركّب
        \emph{مبرهنة كوروفكين}: $L_nf \to f$ بانتظام من أجل
        \emph{كل} $f \in C(\intcc{0}{1})$.
  \item تحقق من أن $(B_n)$ يحقق فرضيات كوروفكين:
        فايرشتراس للمرة الثالثة، من ثلاث وحيدات الحدّ.
  \item ليكن $I_n$ مؤثر الاستكمال الأفيني على قطع
        عند العقد $\frac kn$. بيّن أن $I_n$
        خطي موجب، وأن $I_ne_0 = e_0$، $I_ne_1 = e_1$، و
        $\norm{I_ne_2 - e_2}_\infty = \frac{1}{4n^2}$
        \emph{(فعلى كل خلية يكون خطأ الاستكمال الأفيني
        للدالة $t^2$ هو $(t - a)(b - t)$)}. واختم بكوروفكين:
        فالمستكمِلات المضلّعية تتقارب بانتظام من أجل كل
        دالة متصلة $f$.
\end{enumerate}

\medskip
\noindent\textbf{الجزء الرابع --- الأرباح: الكثافة والعزوم
والمشتقات.}
\begin{enumerate}[resume]
  \item بيّن أن كثيرات الحدود ذوات المعاملات \emph{الناطقة}
        كثيفة في $\bigl(C(\intcc{0}{1}),
        \norm\cdot_\infty\bigr)$: فهذا الفضاء الباناخي
        قابل للفصل.
  \item (العزوم تحدّد الدالة) لتكن $f \in
        C(\intcc{0}{1})$ مع $\int_0^1 f(t)\,t^n \dd t = 0$
        من أجل كل $n \in \N$. بيّن أن $\int_0^1 f P = 0$ من أجل كل
        كثير حدود، ثم أن $\int_0^1 f^2 = 0$، ثم أن $f = 0$.
  \item برهن على متطابقة المشتق
        \[
        (B_nf)'(x) = n\sum_{k=0}^{n-1}\Bigl(
        f\Bigl(\frac{k+1}{n}\Bigr) -
        f\Bigl(\frac kn\Bigr)\Bigr)\,
        \binom{n-1}{k}x^k(1-x)^{n-1-k}
        \]
        \emph{(اشتق $p_k$ وأعد الترقيم --- وهو \omterm{thm:b2:series:abel}{جمع
        آبل})}.
  \item لنفترض أن $f$ من الصنف $C^1$. وباستعمال مبرهنة التزايدات المنتهية في
        كل تزايد وبالمقارنة مع $B_{n-1}(f')$، بيّن أن
        $(B_nf)' \to f'$ بانتظام على $\intcc{0}{1}$. واستنتج:
        من أجل $f \in C^1$ توجد كثيرات حدود تتقارب إلى $f$
        \emph{مع} مشتقاتها.
  \item لنفترض أن $f$ من الصنف $C^2$. وبتايلور--لاغرانج عند $x$ بيّن أن
        \[
        \abs{B_nf(x) - f(x)} \leq
        \frac{\norm{f''}_\infty}{2}\cdot\frac{x(1-x)}{n}
        \leq \frac{\norm{f''}_\infty}{8n} :
        \]
        فترقّي الملاسةُ المعدلَ من $n^{-1/2}$ إلى
        $n^{-1}$.
\end{enumerate}

\medskip
\noindent\textbf{الجزء الخامس --- التشبّع: مبرهنة
فورونوفسكايا.}
\begin{enumerate}[resume]
  \item برهن على متطابقة العزم الرابع
        \[
        \sum_k (k - nx)^4 p_k(x)
        = nx(1-x)\bigl(1 + 3(n-2)x(1-x)\bigr) \leq n^2
        \quad (n \geq 1)
        \]
        \emph{(انشر $k^4$ بالعامليات النازلة $k(k-1)\cdots$
        واستعمل حيلة الاشتقاق في السؤال 2 مرتين
        أخريين)}.
  \item (فورونوفسكايا) لتكن $f$ من الصنف $C^2$ وليكن $x \in
        \intcc{0}{1}$. وبكتابة $f(t) = f(x) + f'(x)(t-x) +
        \frac{f''(x)}2(t-x)^2 + \eta(t)(t-x)^2$ مع $\eta$
        محدودة و$\eta(t) \to 0$ حين $t \to x$، برهن على أن
        \[
        n\bigl(B_nf(x) - f(x)\bigr)
        \xrightarrow[n\to\infty]{}
        \frac{x(1-x)}{2}\,f''(x)
        \]
        \emph{(افصل مجموع $\eta$ عند $\abs{t - x} \leq
        \delta$؛ وتحكّم في الجزء البعيد بالسؤال 23)}. ومن ثم
        فإن خطأ السؤال 22 دقيق في الرتبة \emph{وفي}
        الثابت: أي إن $B_n$ \emph{يتشبّع} عند $\frac1n$، مهما كانت
        $f$ ملساء --- وقارِن بالنتيجة
        \cref{exo:b2:funcseq:7}.
  \item تركيب. بجملة واحدة لكل بند: (1) ماذا اشترت الإيجابية
        وحدها (الجزآن الأول والثالث)؛ (2) وأين دخل تراص
        الفترة $\intcc{0}{1}$ في كل جزء؛ (3) ولماذا تكفي ثلاث
        دوال اختبار في مبرهنة كوروفكين؛ (4) وما المقايضة التي
        يعقدها برنشتاين (معدل $n^{-1/2}$ متين من أجل $f$
        الخشنة، لكن بسقف $\frac1n$ من أجل $f$ الملساء)، وأي
        فصل من هذا الكتاب سيلعب اللعبة نفسها بكثيرات
        الحدود المثلثية.
\end{enumerate}
\end{problem}
